%% kapa: exam=1 round=19 year=2008 period=1 form=A subject=appraisal_law subject_ko=부동산관계법규 seq=11 number=51
%% source: Q-Net 감정평가사 기출문제 (cst00309) artlSeq=1001962 file=extracted/1008/1001962/2008년 제19회 감정평가사 1차시험 기출문제/1교시.hwp
%% conversion: kordoc 4.24.51 -> tools/gen1.py (auto); review=proofread-sonnet; flags=none
\begin{question}{51}
국토의 계획 및 이용에 관한 법령에 의한 벌칙의 최고액이 가장 경미한 것은?
\begin{choices}
  \item 공동구에 수용하여야 하는 시설을 공동구에 수용하지 아니하는 행위
  \item 허가를 받지 않고 공동구를 점용하는 행위
  \item 사위 그 밖의 부정한 방법으로 개발행위허가를 받는 행위
  \item 정당한 권한을 가진 공무원이 개발행위에 대한 업무상황을 검사하는 것을 거부하거나 방해하는 행위
  \item 국토해양부장관이 필요하다고 인정하여 개발행위허가를 받은 자에 대하여 감독상 필요한 보고를 하게 한 경우 \textbf{그} 보고를 하지 아니한 행위
\end{choices}
\end{question}
\answer{5}
